\documentclass{article}
\usepackage{amsmath}
\newcommand{\argmin}{\arg\!\min} 
\newcommand{\argmax}{\arg\!\max} 
\usepackage{amssymb}
\usepackage{amsfonts}
\usepackage[T1]{fontenc}
\usepackage{bm}
\usepackage{array}
\usepackage{graphicx}
\usepackage[utf8]{inputenc}

\title{Decentralized admitance control client}
\date{date}
\author{Marko Guberina}

\begin{document}
\maketitle
\paragraph{Setting} There are 2 robots which carry a thing together.
They share the wrenches they read among each other. The control signal 
(target end-effector position and velocity) and is somehow computed somewhere.
The computer that computes the signal runs ROS1. Hence it needs to receive
a ROS1 wrench message, and it sends out some message for the control singal.

One robot, Universal Robots’ UR5e, is a 6-dof collaborative manipulator with a 5kg payload.
Communication with the robot is achieved with via the Real-Time Data Exchange (RTDE) interface,
with ur\_rtde for client-side implementation. The highest available communication frequency is 500Hz.
Wrench sensing is obtained from UR5e’s built-in force-torque sensor via the same interface.
The Robotiq 2F-85 gripper is connected to the robot’s wrist tool flange, and controlled via RS485 serial communication.
The client-side communication protocol with the gripper is implemented in Python.
The CLIK control law is implemented in a custom library written in Python, leveraging Pinocchio’s
implementation of rigid body algorithms.

Due to versioning incompatibilities between the library and ROS Melodic, the computer controlling
the robot (running the control library utilizing ur\_rtde) can not run ROS1.
Hence communication with the ROS stack has to be handled otherwise.
In particular, we have a thin translation layer
between ROS messages and protobuf running in a Docker container with negligible overhead.
This works because protobuf allow the compilation of messages to old standards 
(even Python's pickling changed in the meantime (although specifying an old pickler should work maybe idk)).

In this directory there is an 
\begin{itemize}
	\item SMC-SMC version of the CLIK executor and a command server
giving end-effector positions and velocities,
\item and the programs to run in the thin Docker containter - same CLIK exector runs on "this computer",
	while these files are to be copy pasted (along with compiled messages) to this Docker intermediary
\end{itemize}







\end{document}
